\section*{Hoofdstuk \ref{ch:b1:cross-sections} --- Geologische doorsneden}

\begin{solution}{exo:b1:cross-sections:1}
Op 1:25\,000 is \qty{1}{cm} op de kaart \qty{250}{m} op de grond. De verticale schaal is
\qty{1}{cm} per \qty{100}{m}, dus de verticale schaal is $250/100 = 2.5$ keer de horizontale: $e = 2.5$.
\end{solution}

\begin{solution}{exo:b1:cross-sections:2}
$\gamma = \ang{90} - \ang{60} = \ang{30}$, en $\tan\beta = \tan\ang{40}
\cos\ang{30} = 0.8391\times0.8660 = 0.7267$: $\beta = \ang{36.0}$, de laag daalt naar het oostelijke
uiteinde van de doorsnede. In de doorsnede met \omterm{def:b1:spherical-astronomy:horizontal}{azimut} \ang{000} is $\gamma = \ang{90}$ en $\beta = 0$: de
doorsnede loopt langs de \omterm{def:b1:geological-maps:attitude}{strijk} en de laag verschijnt horizontaal.
\end{solution}

\begin{solution}{exo:b1:cross-sections:3}
$t = h\cos\delta = 120\times\cos\ang{12.16} = 120\times0.9776 =
\qty{117.3}{m}$.
\end{solution}

\begin{solution}{exo:b1:cross-sections:4}
$\tan\ang{12} = 0.2126$. Met $e = 2$: $\arctan 0.4251 = \ang{23.0}$; met
$e = 3$: $\arctan 0.6377 = \ang{32.5}$; met $e = 5$: $\arctan 1.0628 =
\ang{46.7}$. De hoek groeit veel minder dan $e$ zelf, en bereikt nooit
\ang{90}.
\end{solution}

\begin{solution}{exo:b1:cross-sections:5}
(a) $t = 300\sin\ang{25} = \qty{126.8}{m}$. (b) $t = 300\sin\ang{35} =
\qty{172.1}{m}$. (c) $t = 300\sin\ang{15} = \qty{77.6}{m}$. Dezelfde ontsluitingsbreedte
kan diktes betekenen die meer dan een factor twee verschillen: de helling van het terrein
moet altijd in rekening worden gebracht.
\end{solution}

\begin{solution}{exo:b1:cross-sections:6}
Met \omterm{def:b1:geological-maps:attitude}{hellingrichting} $\alpha$ en helling $\delta$ geeft de stelling in de noord--zuid doorsnede
($\theta = 0$) $\tan\ang{20} = \tan\delta\cos\alpha$, en in de oost--west doorsnede
($\theta = \ang{90}$) $\tan\ang{30} =
\tan\delta\cos(\alpha - \ang{90}) = \tan\delta\sin\alpha$. Delen geeft
$\tan\alpha = 0.5774/0.3640 = 1.586$, $\alpha = \ang{057.8}$; de kwadraten optellen geeft
$\tan\delta = \sqrt{0.3640^2 + 0.5774^2} = 0.6825$, $\delta =
\ang{34.3}$. Beide \omterm{def:b1:cross-sections:apparent-dip}{schijnbare hellingen} zijn kleiner dan de ware helling, zoals het moet zijn.
\end{solution}

\begin{solution}{exo:b1:cross-sections:7}
$100/\tan\ang{12.16} = 100/0.2155 = \qty{464.1}{m}$. Op de kaart liggen de \omterm{def:b1:cross-sections:structure-contour}{structuurlijnen} van de
oostflank inderdaad iets minder dan een halve kilometer uit elkaar, gelijk verdeeld en evenwijdig: de
poot is een vlak.
\end{solution}

\begin{solution}{exo:b1:cross-sections:8}
B1: $524.02 - 488.98 = \qty{35.04}{m}$; B3: $529.98 - 494.93 =
\qty{35.05}{m}$. Boven B2 is de basis van de schalie de top van de kalksteen, \qty{180}{m} boven haar
basis: $280.16 + 180 = \qty{460.16}{m}$, hoger dan het maaiveld (\qty{380.64}{m}), wat de erosie er
verklaart. Tussen B1 en B2 stijgt het contact met $460.16 - 35.04 = \qty{425.12}{m}$
over \qty{2000}{m}: een helling van $0.2126$, een \omterm{def:b1:cross-sections:apparent-dip}{schijnbare helling} van \ang{12.0}. B1
en B3, op gelijke afstand aan weerszijden van B2, vinden elk contact op dezelfde hoogte binnen
\qty{1}{cm}: de \omterm{def:b1:geological-maps:fold}{plooi} is symmetrisch om een verticaal assenvlak bij $x = \qty{3.5}{km}$, haar
beide poten helden even sterk, en geen breuk snijdt de doorsnede tussen hen in.
\end{solution}

\begin{solution}{exo:b1:cross-sections:9}
\omterm{def:b1:cross-sections:throw}{Horizontale verplaatsing} $= 150/\tan\ang{50} = 150/1.1918 = \qty{125.9}{m}$. In de doorsnede met
\omterm{def:b1:spherical-astronomy:horizontal}{azimut} \ang{150} is $\gamma = \ang{30}$ en $\tan\beta = 1.1918\times0.8660 =
1.0321$: de breuk verschijnt met de \omterm{def:b1:cross-sections:apparent-dip}{schijnbare helling} \ang{45.9}. De afsnijlijnen van een horizontale laag liggen
op dat spoor \qty{150}{m} verticaal uit elkaar, dus $150/1.0321 = \qty{145.3}{m}$ langs de doorsnede,
meer dan de \omterm{def:b1:cross-sections:throw}{horizontale verplaatsing}, omdat de doorsnede de breuk schuin snijdt.
\end{solution}

\begin{solution}{exo:b1:cross-sections:10}
$\cos\gamma = \tan\ang{20}/\tan\ang{40} = 0.3640/0.8391 = 0.4338$, dus
$\gamma = \ang{64.3}$: de doorsneden met \omterm{def:b1:spherical-astronomy:horizontal}{azimut} $\alpha \pm \ang{64.3}$ (en de tegengestelde richtingen, waar de
laag stijgt). De \omterm{def:b1:cross-sections:apparent-dip}{schijnbare helling} is \ang{40.0}, \ang{39.6}, \ang{38.3} en \ang{36.0} voor $\gamma = 0$, \ang{10},
\ang{20}, \ang{30}, maar \ang{22.8}, \ang{16.0}, \ang{8.3} en $0$ voor
$\gamma = \ang{60}$, \ang{70}, \ang{80}, \ang{90}. Inderdaad, door $\tan\beta = \tan\delta\cos\gamma$ af te
leiden,
$\dfrac{\mathrm d\beta}{\mathrm d\gamma} =
-\dfrac{\tan\delta\sin\gamma}{1 + \tan^2\delta\cos^2\gamma}$, wat nul is bij $\gamma = 0$ en gelijk is aan
$-\tan\delta$ bij $\gamma = \ang{90}$: een doorsnede die weinig van de \omterm{def:b1:geological-maps:attitude}{hellingrichting} afwijkt, toont bijna de
ware helling, terwijl nabij de \omterm{def:b1:geological-maps:attitude}{strijk} elke graad \omterm{def:b1:spherical-astronomy:horizontal}{azimut} de \omterm{def:b1:cross-sections:apparent-dip}{schijnbare helling} met bijna een graad verandert.
\end{solution}

\begin{solution}{exo:b1:cross-sections:11}
De lagen helen naar het oosten, dus de normaal op de gelaagdheid, naar beneden gevolgd, helt \ang{12.16} uit de
verticale naar het westen, naar \ang{260.7}. (a) Een verticale boring maakt de hoek $\psi = \ang{12.16}$ ermee:
$176.0/\cos\ang{12.16} = \qty{180.0}{m}$, de \omterm{def:b1:cross-sections:thickness}{verticale dikte}. (b) Een boring die \ang{60} naar \ang{260.7} helt, is
\ang{30} uit de verticale gekanteld, aan dezelfde kant als de normaal: $\psi = \ang{30} - \ang{12.16} =
\ang{17.84}$, $\cos\psi = 0.9519$, en de kern is $176.0/0.9519 =
\qty{184.9}{m}$ lang. (c) Naar \ang{080.7} is zij aan de andere kant gekanteld: $\psi = \ang{30} + \ang{12.16} = \ang{42.16}$,
$\cos\psi = 0.7413$, en de kern is $176.0/0.7413 = \qty{237.4}{m}$ lang. De kern is het kortst, gelijk aan
$t$, wanneer de boring langs de normaal loopt, dat wil zeggen \ang{77.84} helt naar \ang{260.7}.
\end{solution}

\begin{solution}{exo:b1:cross-sections:12}
In het oost--west verticale vlak, met de hoeken gemeten vanaf het oosten, stijgen de bovenste lagen naar het oosten
onder \ang{20} en dalen de onderste lagen naar het oosten onder \ang{70}. Door de \omterm{def:b1:geological-maps:geological-map}{ontsluiting} om de \omterm{def:b1:geological-maps:attitude}{strijk} te
draaien met \ang{20}, oostzijde omlaag, worden de bovenste lagen horizontaal en komen de onderste op
$\ang{70} + \ang{20} = \ang{90}$: zij stonden verticaal toen de bovenste lagen werden afgezet. Geschiedenis, van de
oudste gebeurtenis: de onderste lagen worden horizontaal afgezet; zij worden op hun kant gekanteld, tot verticaal,
en opgeheven; hun randen worden tot een oppervlak afgesleten; de bovenste lagen worden over dat oppervlak afgezet;
het geheel wordt \ang{20} naar het westen gekanteld; erosie snijdt het huidige maaiveld. Het is de geschiedenis van
Huttons \omterm{def:b1:geological-maps:geological-map}{ontsluiting} bij Siccar Point.
\end{solution}

\begin{solution}{pb:b1:cross-sections:1}
\textbf{1.} De dikte boven de tunnel is $z - 200$: het grootst bij
$x = \qty{1}{km}$, $556 - 200 = \qty{356}{m}$; het kleinst bij $x =
\qty{9}{km}$, $339 - 200 = \qty{139}{m}$.

\textbf{2.} \qty{1}{km} hoogte meet \qty{4.2}{cm} en \qty{1}{km} afstand
\qty{1.4}{cm}: $e = 4.2/1.4 = 3$.

\textbf{3.} $380.64 - 220.48 = \qty{160.16}{m}$.

\textbf{4.} In B2 ligt de basis van de kalksteen op $380.64 - 100.48 =
\qty{280.16}{m}$ en de basis van de zandsteen op \qty{160.16}{m}: de
tunnel, op \qty{200}{m}, zit in de zandsteen.

\textbf{5.} B1: $524.02 - 668.98 = \qty{-144.96}{m}$; B3: $529.98 - 674.93 =
\qty{-144.95}{m}$.

\textbf{6.} Van west naar oost: krijt, schalie, kalksteen, zandsteen (in de kern van de \omterm{def:b1:geological-maps:fold}{antiklinaal}),
kalksteen, schalie, en opnieuw krijt tot $x = \qty{9}{km}$, door de \omterm{def:b1:geological-maps:fold}{synklinaal} en over de breuk heen.

\textbf{7.} $\gamma = \ang{80.7} - \ang{90} = \ang{-9.3}$; $\tan\beta =
\tan\ang{12.16}\cos\ang{9.3} = 0.2155\times0.9869 = 0.2126$: $\beta =
\ang{12.0}$, naar het oosten dalend.

\textbf{8.} $(280.16 + 144.95)/2000 = 425.11/2000 = 0.2126$: \ang{12.0}.

\textbf{9.} $\arctan(3\times0.2126) = \arctan 0.6377 = \ang{32.5}$.

\textbf{10.} $\gamma = \ang{120} - \ang{90} = \ang{30}$; $\tan\beta =
\tan\ang{60}\cos\ang{30} = 1.732\times0.866 = 1.500$: $\beta = \ang{56.3}$.

\textbf{11.} De poten helen naar \ang{080.7} of \ang{279.3}, slechts \ang{9.3} van de lijn van de
doorsnede, en $\cos\ang{9.3} = 0.987$; de breuk helt naar \ang{120}, \ang{30} daarvan, en
$\cos\ang{30} = 0.866$.

\textbf{12.} Het breukvlak is $z = 250 - \tan\ang{60}\,[(x - 4300)
\sin\ang{120} + (y - 1000)\cos\ang{120}]$. Met $y = 6000$ is $(y - 1000)
\cos\ang{120} = -2500$, dus $(x - 4300)\sin\ang{120} = 2500 - (z - 250)/
\tan\ang{60}$, en $1/\tan\ang{60} = \tan\ang{30}$.

\textbf{13.} Op $z = 200$: $x = 4300 + (2500 + 50\times0.5774)/0.8660 =
4300 + 2528.9/0.8660 = \qty{7220}{m}$, dat is \qty{7.22}{km} van C.

\textbf{14.} \omterm{def:b1:cross-sections:throw}{Horizontale verplaatsing} $= 200/\tan\ang{60} = \qty{115.5}{m}$; afstand langs
C--D $= 200/1.500 = \qty{133.3}{m}$.

\textbf{15.} De breuk helt oost-zuidoostwaarts en is normaal: het hangende blok, naar het oosten, is omlaag bewogen.
Op tunnelhoogte ligt aan beide kanten van de breuk krijt: ten westen ervan het krijt van de \omterm{def:b1:geological-maps:fold}{synklinaal}; ten oosten
ervan krijt dat \qty{200}{m} omlaag is gebracht.

\textbf{16.} De tunnel blijft over de breuk in krijt, dus verandert het gesteente aan het front niet en kondigt niets
de breuk aan; alleen een zone van verbrijzeld gesteente, en misschien water, doet dat. De ingenieurs voorspellen de
plaats ervan uit de doorsnede (hier \qty{7.22}{km}), boren proefgaten vooruit van de machine als zij nadert, en bereiden
steunwerk voor de verbrijzelde zone voor.

\textbf{17.} $668.98 - 488.98 = \qty{180.00}{m}$.

\textbf{18.} $t = 180\cos\ang{12.16} = \qty{176.0}{m}$.

\textbf{19.} Op de kruin ($x = 3500$) ligt de basis van de zandsteen op \qty{160.16}{m} (vraag 3); zij daalt aan
beide poten met de \omterm{def:b1:cross-sections:apparent-dip}{schijnbare helling} \ang{12.0} van vraag 7, dat wil zeggen met $\tan\ang{12.0}$ per meter horizontale
afstand van de as.

\textbf{20.} De basis van de kalksteen ligt \qty{120}{m} boven de basis van de zandsteen, op $280.16 - u\tan\ang{12.0}$ met
$u = |x - 3500|$, en haar top \qty{180}{m} hoger, op $460.16 - u\tan\ang{12.0}$. Met $\tan\ang{12.0} =
0.21256$: de basis ligt onder \qty{200}{m} wanneer $u > 80.16/0.21256 = 377.1$, en de top ligt boven
\qty{200}{m} wanneer $u < 260.16/0.21256 = 1224.0$.

\textbf{21.} $x = 3500 \pm 377.1$ en $3500 \pm 1224.0$: de tunnel zit in de kalksteen van
\qty{2276}{m} tot \qty{3123}{m} en van \qty{3877}{m} tot \qty{4724}{m}.

\textbf{22.} Langs een horizontale lijn worden de basis en de top van de kalksteen, \qty{180}{m} verticaal uit elkaar en
beide dalend met $\tan\beta$, gekruist op een horizontale afstand $h/\tan\beta$ van elkaar: $180/0.21256 = \qty{846.8}{m}$
op elke poot, hetzelfde tot op afronding als $1224.0 - 377.1 = 846.9$.

\textbf{23.} $2\times846.8 = \qty{1693.7}{m}$: de tunnel boort \textbf{ongeveer
\qty[text-series-to-math=true]{1694}{m} kalksteen}, meer dan een vijfde van zijn
\qty{8}{km} lengte, in twee stukken van \qty{847}{m} aan weerszijden van de zandsteenkern van de \omterm{def:b1:geological-maps:fold}{antiklinaal}.
\end{solution}
